\documentclass[aspectratio=169,10pt]{beamer}
\usetheme[numbering=fraction,progressbar=frametitle]{metropolis}
\usepackage{fontspec}
\setsansfont{Latin Modern Sans}
\usepackage{amsmath,amssymb,mathtools,graphicx,booktabs}
\usepackage{polyglossia}
\setdefaultlanguage{english}
\setbeamertemplate{navigation symbols}{}
\setbeamerfont{frametitle}{size=\large}
\setbeamerfont{normal text}{size=\normalsize}
\setlength{\parskip}{.45em}
\hypersetup{pdfauthor={},pdfsubject={Lecture slides},pdfkeywords={},pdftitle={Term review and consolidation}}
\providecommand{\N}{\mathbb{N}}
\providecommand{\Z}{\mathbb{Z}}
\providecommand{\D}{\mathbb{D}}
\providecommand{\Q}{\mathbb{Q}}
\providecommand{\R}{\mathbb{R}}
\providecommand{\C}{\mathbb{C}}
\providecommand{\vect}[1]{\overrightarrow{#1}}
\providecommand{\vi}{\vec{\imath}}
\providecommand{\vj}{\vec{\jmath}}
\providecommand{\vk}{\vec{k}}
\providecommand{\norm}[1]{\left\lVert #1\right\rVert}
\providecommand{\vc}[2]{\begin{pmatrix}#1\\#2\end{pmatrix}}
\providecommand{\deter}[2]{\det\left(#1,#2\right)}
\providecommand{\abs}[1]{\left\lvert #1\right\rvert}
\providecommand{\dd}{\,\mathrm{d}}
\providecommand{\e}{\mathrm{e}}
\providecommand{\Card}{\operatorname{Card}}
\providecommand{\E}{\operatorname{E}}
\providecommand{\V}{\operatorname{V}}
\providecommand{\degre}{^{\circ}}
\providecommand{\un}[1]{\,\mathrm{#1}}
\providecommand{\sci}[2]{#1\times10^{#2}}
\providecommand{\pgcd}{\operatorname{PGCD}}
\providecommand{\ppcm}{\operatorname{PPCM}}
\title{Term review and consolidation}
\subtitle{Lecture and worked examples}
\date{}
\author{}
\begin{document}
\begin{frame}[plain]
\vfill
{\usebeamerfont{title}\inserttitle\par}
\medskip
{\usebeamerfont{subtitle}\insertsubtitle\par}
\medskip
{\usebeamercolor[fg]{progress bar}\rule{\textwidth}{.5pt}}
\vfill
\end{frame}
\begin{frame}{Individual review}
Work without the course notes for 25 minutes, then check with the methods in the booklet. For each question, note: solved alone, solved with help, or to redo. Justify the answers.
\end{frame}
\begin{frame}{Individual review}
\begin{enumerate}
\setcounter{enumi}{0}
\item Reduce \(90/126\) to lowest terms. Prove that the sum of two odd integers is even.

\item With \(A(-1,2)\) and \(B(3,-2)\), calculate \(\vect{AB}\), \(AB\) and the midpoint of \([AB]\) in an orthonormal frame.

\item Determine an equation of \((AB)\); does \(C(5,-4)\) belong to this line?

\item Calculate \([-2,4]\cap(1,7]\) and \([-2,4]\cup(1,7]\).

\item Solve \(|x-1|\leq4\) and explain why \(1/3\notin\D\).

\end{enumerate}
\end{frame}
\begin{frame}{Worked answer: Individual review}
1. \(\operatorname{gcd}(90,126)=18\), so \(90/126=5/7\). If \(u=2k+1\) and \(v=2\ell+1\), then \(u+v=2(k+\ell+1)\) is even.

2. \(\vect{AB}=(4,-4)\); \(AB=4\sqrt2\); midpoint \(I(1,0)\).

3. \(y=-x+1\). For \(C\), \(-4=-5+1\): it belongs to \((AB)\).

4. Intersection \((1,4]\); union \([-2,7]\).

5. \(x\in[-3,5]\). Writing \(1/3=a/10^n\) would require \(10^n\) to be divisible by \(3\), which is impossible.
\end{frame}
\begin{frame}{Correcting an error and revisiting a method}
Correct: (a) “\(\vect{AB}=(x_A-x_B,y_A-y_B)\)”; (b) “The zero vector directs every line”; (c) “\(|x-2|\leq3\) gives \(x\in[2,3]\)”; (d) “\(1/3\approx0.333\), so \(1/3\) is decimal”.
\end{frame}
\begin{frame}{Worked answer: Correcting an error and revisiting a method}
(a) Subtract the origin from the endpoint: \((x_B-x_A,y_B-y_A)\). (b) A direction vector must be nonzero. (c) The distance to \(2\) is at most \(3\): \(x\in[-1,5]\). (d) \(0.333\) differs from \(1/3\): \(3\times0.333=0.999\ne1\). A decimal approximation does not determine the nature of the approximated number.
\end{frame}
\begin{frame}{Correcting an error and revisiting a method}
\begin{enumerate}
\setcounter{enumi}{0}
\item Reduce \(72/108\), then \(\dfrac34-\dfrac16\).

\item With \(P(1,1)\) and \(Q(5,3)\), determine the midpoint, a direction vector and an equation of \((PQ)\).

\item Determine \((-\infty,3]\cap[-1,+\infty)\), then solve \(|x+2|<1\).

\end{enumerate}
\end{frame}
\begin{frame}{Worked answer: Correcting an error and revisiting a method}
1. \(2/3\); \(7/12\). 2. Midpoint \((3,2)\), direction vector \((2,1)\), equation \(x-2y+1=0\). 3. Intersection \([-1,3]\); solution of the inequality \((-3,-1)\).
\end{frame}
\begin{frame}{Applying and self-assessing}
Two displacements are represented by \(\vec u(6,4)\) and \(\vec v(9,6)\).

\begin{enumerate}
\setcounter{enumi}{0}\item Show that they are collinear and find the coefficient relating \(\vec v\) to \(\vec u\).

\item The starting point is \(A(-2,1)\). Determine the endpoints \(B\) and \(C\) associated with the displacements \(\vec u\) and \(\vec v\).

\item Find an equation of the line through these three points.

\item A point \(M(x,y)\) of this line must satisfy \(|x-4|\leq3\). Give the interval of possible abscissae.

\end{enumerate}
\end{frame}
\begin{frame}{Worked answer: Applying and self-assessing}
1. \(\vec v=\dfrac32\vec u\). 2. \(B(4,5)\) and \(C(7,7)\). 3. A direction vector is \((3,2)\); the equation is \(2x-3y+7=0\), satisfied by the three points. 4. \(|x-4|\leq3\iff x\in[1,7]\).
\end{frame}
\begin{frame}{Applying and self-assessing}
My priority:  

The mistake to avoid:  

The method to reuse:  

An exercise to redo without help:
\end{frame}
\end{document}
